\section{Introduction}
\label{sec:intro}

\paragraph{Abstract.}
Vision-Language Models (\vlm{}s) routinely accept user-supplied images alongside benign text prompts. We ask whether a single image, perceptually indistinguishable from a real photo, can carry hidden instructions that the \vlm{} obeys at inference time. Building on two existing papers --- Universal Adversarial Attack on Aligned Multimodal LLMs \citep{rahmatullaev2025universal} and AnyAttack \citep{zhang2025anyattack} --- we compose a three-stage pipeline (\textsc{VisInject}) that turns one universal adversarial image into a stream of perceptually invisible adversarial photos at PSNR $\approx 25.2$~dB, and we evaluate $6{,}615$ (clean, adversarial) response pairs along two independent axes: \emph{Output Affected} (was the answer disturbed?) and \emph{Target Injected} (did the chosen payload appear?). The dual-axis split shows that broad disruption ($\approx 66\%$) coexists with rare payload delivery ($0.227\%$), and that the few injections that do land cluster on screenshot-like images whose semantics already invite text transcription.

\paragraph{What this report does.}
This is a course final project; we keep the writing oriented toward analysis rather than novelty. The technical contribution is mostly compositional: we build on two existing papers and add a simple but useful evaluation method. The empirical contribution is a careful breakdown of \emph{where injection actually lands}: we show the full $7 \times 7 \times 3$ design, run it end-to-end, and analyse the gap between disruption and injection that earlier judge-based evaluations had hidden.

\paragraph{Headline takeaways.}
\begin{itemize}
  \item Disruption is easy and broad: every Qwen / DeepSeek model in our set is disturbed on $\geq 98\%$ of pairs at $\Linf = 16/255$.
  \item Injection is rare and clustered: $0.227\%$ overall, with the two confirmed cases both being URL injections on a code screenshot.
  \item Architecture beats parameter count: BLIP-2 is fully immune across $2{,}205$ pairs, while smaller transformer-style \vlm{}s are not.
  \item No transfer to GPT-4o: the strongest small-model injection fails on the frontier model in a manual transfer test.
  \item As a side effect, the released HuggingFace dataset attracts $\sim$300 downloads in its first month, suggesting the pipeline + curated cases are usable beyond this report.
\end{itemize}

\paragraph{Structure.}
\S\ref{sec:threat} states the threat model. \S\ref{sec:method} summarises the two papers we build on and our dual-dim evaluation. \S\ref{sec:experiments} is the longest section: it lays out the full experiment design (every prompt, every test image, every VLM ensemble). \S\ref{sec:results}--\S\ref{sec:cases} cover aggregate results and case studies. \S\ref{sec:discussion} discusses why the gap between disruption and injection is what it is. \S\ref{sec:future} sketches what \textsc{VisInject} v2 would look like.
